\section{Methods: Travel Time and Market Analysis}
\label{sec:methods}

\subsection{Bus Travel Time Model}

Scheduled bus travel time on a T1 managed lane corridor is modeled as:

\begin{equation}
  T_{\text{bus}} = \frac{d}{v_{\text{bus}}} + N_{\text{stops}} \cdot t_{\text{stop}}
  \label{eq:travel-time}
\end{equation}

where $d$ is route distance in miles, $v_{\text{bus}}$ is the effective average cruising speed in mph, $N_{\text{stops}}$ is the number of intermediate stops, and $t_{\text{stop}}$ is the dwell time per stop in minutes.

\paragraph{Effective cruising speed.} T1 managed lanes sustain 65 mph under design conditions (PTI $\leq$ 1.15). We apply a 3 mph deduction for stop-approach deceleration and acceleration cycles integrated over the full route, yielding $v_{\text{bus}} = 62$ mph. This is conservative relative to express coach performance on uncongested European motorways, where 65--68 mph averages are common \citep{FlixBus2023}.

\paragraph{Stop dwell time.} Each intermediate stop (T1/T2 hub platform boarding) is modeled at $t_{\text{stop}} = 8$ minutes: 3 minutes deceleration and platform approach, 3 minutes passenger boarding and alighting with wheelchair-accessible coach protocol, 2 minutes departure and acceleration to cruising speed. Express routes with 3--4 intermediate stops therefore add 24--32 minutes to the free-flow equivalent. We model all 12 corridors with 3 intermediate stops as a baseline; Section~\ref{sec:travel-time} notes where stop counts differ by corridor.

\paragraph{Stop-Penalty Adjustment.}
The travel time model incorporates a stop-penalty for intermediate T1/T2 hub stops.
Each intermediate stop (boarding, alighting, and acceleration/deceleration cycle)
adds approximately 8 minutes of elapsed time:

\begin{equation}
T_{\text{bus}} = \frac{d}{v_{\text{cruising}}} + N_{\text{stops}} \times t_{\text{stop}}
\label{eq:bus-time}
\end{equation}

where $d$ is distance, $v_{\text{cruising}} = 65$~mph (managed lane design speed),
$N_{\text{stops}}$ is the number of intermediate T1/T2 hub stops, and $t_{\text{stop}}
= 8$~min per stop. The number of intermediate stops is estimated as:

\begin{equation}
N_{\text{stops}} = \left\lfloor \frac{d}{150} \right\rfloor
\end{equation}

reflecting T1/T2 hub spacing of approximately one stop per 150 miles along T1 corridors.
The resulting effective average speed $\bar{v}$ (total distance divided by total elapsed
time) is lower than the cruising speed:

\begin{align}
\bar{v} &= \frac{d}{d/v_{\text{cruising}} + N_{\text{stops}} \times t_{\text{stop}}/60} \\
        &\approx 59\text{--}61\text{ mph for corridors } 300\text{--}900\text{ miles}
\end{align}

This effective speed of 59--61 mph still represents a substantial advantage over current
Greyhound operations (mean effective speed approximately 44 mph including stops and layovers)
but is more conservative than a no-stop managed-lane assumption. The effective speed of
59--61 mph is consistent with the $v_{\text{bus}} = 62$~mph operating speed used in
Equation~\ref{eq:travel-time}: the 3-event stop-penalty at 8 minutes per stop over
corridors of 300--900 miles reduces the arithmetic mean speed by 1--3 mph relative to
the cruising speed. Travel time values in the 12-corridor comparison table
(Section~\ref{sec:travel-time}) reflect this stop-penalty model using $v_{\text{bus}} = 62$~mph
as the operating speed (not 65 mph; the 3 mph deduction for stop-approach cycles is already
incorporated into $v_{\text{bus}}$).

\paragraph{PTI buffer and schedule reliability.} The scheduled time posted to passengers is derived from Equation~\ref{eq:travel-time} as the \emph{median-day} travel time. Under PTI 1.15, the 95th percentile travel time is $T_{\text{bus}} \times 1.15$, meaning the scheduled time is reliable for 95\% of trips with no additional buffer needed. Operators have the option to add a further 5--10 minute terminal arrival buffer to achieve 98\%+ on-time performance at published arrival times; this is not incorporated in the travel time comparisons below, which use the unpadded median-day time.

\paragraph{Comparison baselines.} Current Greyhound and FlixBus travel times are taken from published operator schedules (winter 2025--2026). Where Greyhound does not operate (service withdrawn), we estimate current-network travel time from alternative connections (minimum-transfer composite routes). Amtrak travel times are from published FY2023 schedules \citep{Amtrak2023}; we use the scheduled (not historical on-time) time to give Amtrak the most favorable comparison. Drive times are Google Maps free-flow (no congestion). Air times include 1.5 hours airport processing each end plus scheduled flight duration from Bureau of Transportation Statistics O-D data \citep{BTS_intermodal2023}.

\subsection{PTI Translation for Current-Network Baselines}

Current Greyhound schedules on I-80 already incorporate de facto PTI buffers: operators know they cannot run reliably and pad schedules accordingly. The 18--22 hour published Chicago--New York Greyhound time is not a 62 mph average; it is a 35--43 mph average that includes 3--4 hours of schedule padding absorbing PTI 1.86 unreliability. We treat published operator times as the relevant comparison because that is the service experience available to travelers, regardless of what free-flow speed would theoretically permit.

\subsection{Ridership Estimation}

We estimate corridor ridership using a gravity model calibrated on existing Greyhound/FlixBus routes where ridership data is available \citep{BTS_intermodal2023}:

\begin{equation}
  Q_{ij} = G \cdot P_i^\alpha \cdot P_j^\beta \cdot f_{ij}^\gamma \cdot T_{ij}^\delta
  \label{eq:gravity}
\end{equation}

where $Q_{ij}$ is annual passengers between markets $i$ and $j$, $P_i$ and $P_j$ are metropolitan statistical area populations, $f_{ij}$ is the one-way fare, and $T_{ij}$ is travel time in hours. Parameters $\alpha = 0.85$, $\beta = 0.85$, $\gamma = -1.2$, and $\delta = -0.7$ are calibrated on 18 existing intercity bus corridors for which ridership counts are available in the BTS Intermodal Connectivity Database \citep{BTS_intermodal2023}. The calibration yields $R^2 = 0.79$ in-sample. We apply a 30\% discount to new-market corridors (where no existing bus service operates) to account for awareness effects and the time required to build ridership.

The total market equilibrium estimate of 24 million annual passengers is a steady-state figure (5+ years post-opening); first-year ridership is estimated at 40--55\% of equilibrium for new corridors and 70--80\% for corridors replacing existing services.

\paragraph{Local Access Assumption.}
The ridership gravity model assumes that travelers can reach T1/T1 hub stops. For
general-population travelers with private vehicles, this assumption is reasonable: hub
locations are accessible by car from within a 30-mile radius. For transit-dependent travelers
(zero-vehicle households), the model implicitly assumes local feeder service connects them
to the hub. This feeder assumption is not costed in the operator economics analysis
(Section~\ref{sec:operator-economics}); its omission means the 24 million annual passenger
estimate represents an upper bound for markets where feeder service does not already exist.
Markets served by existing rural transit (FTA 5311 grantees) are partially excepted.

\subsection{Operator Economics Model}

Operating cost per route-day is:

\begin{equation}
  C_{\text{day}} = c_{\text{mile}} \cdot d \cdot N_{\text{buses}} \cdot 2
  \label{eq:cost}
\end{equation}

where $c_{\text{mile}} = \$2.80$ per bus-mile (driver \$0.80, fuel \$1.40, vehicle O\&M \$0.40, insurance \$0.20), $d$ is one-way distance, $N_{\text{buses}}$ is the number of buses per direction per day, and the factor of 2 accounts for the return trip. Revenue is:

\begin{equation}
  R_{\text{day}} = f_{\text{mile}} \cdot d \cdot S \cdot \lambda \cdot N_{\text{buses}} \cdot 2
  \label{eq:revenue}
\end{equation}

where $f_{\text{mile}}$ is the per-mile fare, $S = 50$ seats per coach, and $\lambda$ is the load factor. Break-even load factor is the value of $\lambda$ at which $R_{\text{day}} = C_{\text{day}}$.

The Essential Intercity Transportation (EIT) subsidy analogy follows the structure of the Essential Air Service (EAS) program \citep{EAS_DOT2023}, which subsidizes air service to small communities at an average of \$175/passenger. We propose EIT subsidies capped at \$15/passenger for corridors that are marginally below break-even but serve communities with no other intercity transit option. Section~\ref{sec:economics} applies this model to marginal corridors.
